%%%%%%%%%%%%%%%%%%%%%%%%%%%%%%%%%%%%%%%%%%%%%%%%%%%%
% Artifact Appendix (EuroSys'22 AE template) for
% "How Low Can You Go? The Data-Light SE Challenge"
% Ganguly & Menzies, FSE 2026. arXiv:2512.13524v2
% Repo: https://github.com/KKGanguly/NEO @ 5a0d487 (2026-07-27)
% Drafted 2026-09-14 from repo inspection + smoke test of
% report/plot scripts only. Full optimizer re-runs NOT executed.
%%%%%%%%%%%%%%%%%%%%%%%%%%%%%%%%%%%%%%%%%%%%%%%%%%%%

\appendix
\section{Artifact Appendix}

\subsection{Abstract}
The artifact contains the Python optimizers (LITE, LINE, EZR, RANDOM
wrappers plus DEHB, SMAC and TPE via their public packages), a Lua
scoring harness, precomputed results for all optimizers over the MOOT
benchmark tasks, and scripts that rebuild the paper's summary table
(Table~3, percent-best and $\Delta$ per budget), Figure~4 (percent-best
vs.\ label budget) and Figure~5 (runtime box plots). The cheap path
(table and figures from precomputed results) runs in seconds. The
expensive path (re-running all optimizers on all tasks) is optional.

\subsection{Description \& Requirements}

\subsubsection{How to access}
\url{https://github.com/KKGanguly/NEO}, commit \texttt{5a0d487}
(2026-07-27), MIT license, $\approx$104\,MB cloned (mostly precomputed
results). The benchmark data is \emph{not} in this repository: clone
\url{https://github.com/timm/moot} into \path{NEO/moot/} (the
Makefiles and Lua harness expect \path{NEO/moot/optimize/*/*.csv};
the README's \path{data/moot/} path is wrong).

\subsubsection{Hardware dependencies}
Authors' rig: 4-core 1.30\,GHz Linux, 16\,GB RAM, no GPU. E1--E3 below
run on any laptop. E4 (full re-run) benefits from many cores; the
generated \texttt{commands.sh} launches one process per dataset in
parallel.

\subsubsection{Software dependencies}
Python (README says 3.13; \texttt{.python-version} says 3.11; 3.13
confirmed for E2--E3) with \texttt{requirements.txt} (pandas,
matplotlib, seaborn, scipy, numba, ConfigSpace, dehb, smac, optuna,
optunahub, hebo; CMake/SWIG needed to build SMAC). Lua~5.x, GNU make,
bash, awk, egrep, gawk. Lua is required by \texttt{make comparez}
and is \emph{not} mentioned in the README. No container is provided.

\subsubsection{Benchmarks}
127 multi-objective SE optimization tasks from MOOT (Table~1 of the
paper): software configuration, HPO, product lines, project health,
cost/effort, defect prediction. Note: the committed
\texttt{report.csv} covers 112 tasks; 15 tasks named in paper v2 are
absent from the precomputed results.

\subsection{Set-up}
{\footnotesize\begin{verbatim}
git clone https://github.com/KKGanguly/NEO
cd NEO
git clone https://github.com/timm/moot  # NEO/moot
python3 -m venv .venv
. .venv/bin/activate
pip install -r requirements.txt
brew install lua    # or: apt install lua5.4
\end{verbatim}}
Basic test (no data needed): \texttt{cd experiments/LUA\_run\_all
\&\& make report} must print nothing and rewrite
\path{results/optimization_performance/report.csv}.

\subsection{Evaluation workflow}

\subsubsection{Major Claims}
\begin{itemize}
\item \textit{(C1)} With label budgets of 6--200, simple pool-based
  optimizers (LITE, LINE, RANDOM, EZR) match or beat SMAC, TPE and DEHB
  in percent-best rank and normalized improvement $\Delta$; LITE is best
  overall and DEHB worst. Shown in Table~3 (\S5.5). Reproduced by E1.
\item \textit{(C2)} Percent-best gains flatten after $\approx$100
  labels; small budgets favor LINE/RANDOM, mid budgets favor LITE.
  Shown in Figure~4. Reproduced by E2.
\item \textit{(C3)} RANDOM runs $>$10{,}000$\times$ faster than the
  membership-query optimizers; SMAC is $>$2 orders of magnitude slower
  than LITE/EZR. Shown in Figure~5. Reproduced by E3.
\item \textit{(C4)} The BINGO effect (\S6, Table~4, Figure~7: data
  collapses into few occupied buckets). \textbf{Not covered}: no script
  in the repository regenerates Table~4 or Figure~7.
\end{itemize}

\subsubsection{Experiments}

\medskip\par\noindent\textit{Experiment (E1): Summary table} [5 human-min + 1 compute-sec
from committed scores; Lua rescoring untimed]. Rebuilds Table~3.\\
\textit{[Preparation]}\
Set-up above. \texttt{comparez.out} (Lua
percent-best/$\Delta$ per task) is committed, so the Lua step is
optional.\\
\textit{[Execution]}\
\texttt{cd experiments/LUA\_run\_all; make
report}. To rescore from the raw runs: \texttt{make comparez} first
(runs \texttt{lua run\_all.lua --comparez} once per MOOT csv, reading
\path{results/results_*/} through \texttt{FileResultsReader.py};
RANDOM and LINE are computed inside the Lua script, not from Python
runs).\\
\textit{[Results]}\
\path{results/optimization_performance/report.csv}:
one row per task, columns \texttt{<OPT>-<budget>} and
\texttt{<OPT>-<budget>\_sd}; last row is the column mean. Compare the
mean row with Table~3 (e.g.\ LITE $=$ ACT column: 32, 54, 70, 75, 82,
88, 94; DEHB: 24, 25, 31, 39, 59, 64, 69). Regenerating from the
committed \texttt{comparez.out} gave $\pm$1 differences in a few cells
of the mean row vs.\ the committed csv.


\medskip\par\noindent\textit{Experiment (E2): Figure 4} [2 human-min + 1 compute-sec].\\
\textit{[Execution]}\
\texttt{cd experiments; MPLBACKEND=Agg python3
optim\_performance\_comp.py}. Without \texttt{MPLBACKEND=Agg} the
script blocks on \texttt{plt.show()}.\\
\textit{[Results]}
\path{results/optimization_performance/optimization_performance_comparison.png}:
percent-best vs.\ budget, one line per optimizer; LITE on top from
budget 12 onward, DEHB lowest everywhere (C2).


\medskip\par\noindent\textit{Experiment (E3): Figure 5} [2 human-min + 1 compute-sec].\\
\textit{[Execution]}\
\texttt{cd experiments; MPLBACKEND=Agg python3
performance\_box.py} (reads \texttt{report\_tmp.csv}, written by E1).\\
\textit{[Results]}
\path{results/runtime_plot/families_runtime_boxplot.png}, log-ms
boxes for DEHB, SMAC, TPE, LITE, EZR, LINE, RANDOM; RANDOM lowest,
SMAC highest (C3).


\medskip\par\noindent\textit{Experiment (E4, optional): Re-run optimizers} [10 human-min +
days of compute, not measured]. 127 tasks $\times$ 7 budgets
$\times$ 20 repeats per optimizer.\\
\textit{[Execution]}\
From repo root: \texttt{make generate-commands
NAME=SMAC} (also DEHB, TPE, Active\_Learning, EZR), then
\texttt{cd experiments \&\& ./commands.sh}. Output replaces
\path{results/results_<NAME>/<NAME>/<task>_<budget>.csv} (three
lines: sampled configs, d2h scores, seconds, one entry per repeat).
Then E1--E3. Stochastic; expect Table~3 cells to move by a few points.

\subsection{Notes on Reusability}
New task: drop a MOOT-format csv (goal columns end in \texttt{+}/\texttt{-})
under \path{moot/optimize/<group>/}; \texttt{make generate-commands}
picks it up; add its name to the \texttt{comparez} glob in
\path{experiments/LUA_run_all/Makefile}. New optimizer: subclass
\path{optimizers/base_optimizer.py}, register in
\path{experiments/experiment_runner_cluster.py}, add a results
reader block in \texttt{run\_all.lua}. Budgets and repeats are the
\texttt{--budget}/\texttt{--repeats} flags in the root Makefile.

\subsection{General Notes}
Known gaps an evaluator will hit: (1) README points to
\path{data/moot/}, code expects \path{moot/}; (2) README calls the
summary "Table 4", paper v2 numbers it Table~3; (3) Lua and gawk are
undocumented dependencies; (4) 112 of 127 tasks have precomputed
results; (5) no Docker image; (6) \texttt{plt.show()} in E2 needs a
headless backend on servers.
